\chapter*{Declaration on the Use of Generative Artificial Intelligence}
\addcontentsline{toc}{chapter}{Declaration on the Use of Generative Artificial Intelligence}

\textcolor{eclipsepurple}{\textbf{The text in purple is informational and must be removed from the final version of the document.} The use of generative artificial intelligence must be declared. It is also strongly suggested to state if it was \emph{not} used. Students must read the \href{https://libguides.biblio.usherbrooke.ca/ld.php?content_id=37838948}{guide on the responsible use of generative artificial intelligence in graduate studies}, produced by Myriam Beaudet and Mireille Léger-Rousseau of the Université de Sherbrooke library and archives service, which states that:}

\begin{customquote}
\textcolor{eclipsepurple}{%
\selectlanguage{french}«\,Toute utilisation d'IA ayant contribué au contenu d’une production doit être déclarée. Cette transparence maintient le lien de confiance et permet de se conformer aux obligations de déclaration. La déclaration doit généralement préciser l'usage fait de l'IA et identifier l'outil utilisé ainsi que sa version. Les usages prévus doivent être discutés avec la ou les directions de recherche. […] L'usage non transparent ou non autorisé constitue un manquement à l'intégrité académique et peut entraîner des sanctions selon le Règlement des études (Règlement 2575-009)\,».\selectlanguage{english}%
}
\end{customquote}

\textcolor{eclipsepurple}{The declaration must include the following elements:}

\begin{itemize}
    \item \textbf{Tools used}: [State the name of the tool and its version. E.g.: Claude 3.5 Sonnet (Anthropic), ChatGPT-4 (OpenAI), GitHub Copilot].
    \item \textbf{Nature of the assistance}: [Indicate the types of tasks delegated. E.g.: language editing, Python code generation, French-English translation, generation of visualizations, rewording of sections. See \href{https://libguides.biblio.usherbrooke.ca/ld.php?content_id=37838948}{appendix C of the guide on the responsible use of generative artificial intelligence in graduate studies} if needed].
    \item \textbf{Scope}: [Indicate the chapters or sections concerned. E.g.: Chapter 2 - Literature review, Section 3.2 - Data analysis, figures 4.1 to 4.3.].
    \item \textbf{Type of data provided to the AI}: [Describe the types of content provided AND confirm that no confidential data, including personal information, and no unpublished data or data covered by a confidentiality agreement was transmitted, OR describe the authorizations obtained].
    \item \textbf{Oversight and validation}: [Describe the verification procedures: proofreading, reference validation, testing, replication of analyses, consultation of primary sources, etc.].
    \item \textbf{Safeguards}: [If applicable, describe anonymization procedures, use of local models, etc.].
\end{itemize}
